% Tableau des labs du cours GCI dans NDAP SUKULU (propre au cours).
Le cours compte dix laboratoires, disponibles sous le code \codeUE. Le dernier
réunit les quatre organisations du fil rouge et sert de synthèse.

\begin{tableaucours}{lL{6.2cm}L{1.6cm}L{2.0cm}}{%
    \tableauHeader{Lab} &
    \tableauHeader{Sujet} &
    \tableauHeader{Chapitre} &
    \tableauHeader{Opération}
}
GCI-LAB-01 & Arbitrer le portefeuille de projets de l'ANPP & \ref{chap:alignement}
& MEFOU \\
GCI-LAB-02 & Fixer les capacités cibles de la BANQUE SANAGA &
\ref{chap:cobit-conception} & SANAGA \\
GCI-LAB-03 & Dérouler la cascade des objectifs de la BANQUE SANAGA &
\ref{chap:cobit-conception} & SANAGA \\
GCI-LAB-04 & Évaluer la capacité de BAI06 et de DSS02 & \ref{chap:cobit-capacite}
& SANAGA \\
GCI-LAB-05 & Tester les contrôles généraux informatiques de l'ANPP &
\ref{chap:audit} & MEFOU \\
GCI-LAB-06 & Prioriser et analyser les tickets de KRIBILOG &
\ref{chap:itil-incidents} & KRIBI \\
GCI-LAB-07 & Évaluer et décider un changement de la BANQUE SANAGA &
\ref{chap:itil-changements} & SANAGA \\
GCI-LAB-08 & Disponibilité, pénalités et reprise du suivi des conteneurs &
\ref{chap:itil-niveaux} & KRIBI \\
GCI-LAB-09 & Construire la matrice RACI du modèle de NOUN PAY &
\ref{chap:processus} & NOUN \\
GCI-LAB-10 & Synthèse des quatre cas~: feuilles de route chiffrées &
\ref{chap:cas-integrateurs} & les quatre \\
\end{tableaucours}
\vspace{-0.8em}
\begin{center}{\small\itshape Les laboratoires du cours dans NDAP SUKULU}\end{center}
